\subsection{Retry readiness under the common opportunity budget}
\label{sec:retry-readiness}

We compared immediate and age-four retries at both send budgets under
the seven schedules in Section~\ref{sec:matched-opportunities}. Each
of the 14 pairs used the same protocol source, requests, capture cut,
2,048 transmission opportunities, and cryptographic and durability
interfaces. All 14 fresh immediate-policy runs reproduced their
original rows.

\paragraph{A public retry clock.}
The clock counts the sender's own turns. Each endpoint advances its
counter after recording the opportunity and request, before planning
control traffic. If it sends on turn $k$, an unacknowledged request is
eligible again at $k+4$ at the earliest. An exact retry records a new
emission turn. When fair cyclic arbitration selects the retry class,
the oldest eligible request resumes with its saved ciphertext. Control
planning, cumulative ACK processing, first admission and the other
arbitration rules keep their original code. The sender goes by its
clock and its ACK state. Only the analyst uses delivery history.

% Generated by artifacts/retry-readiness/analyze_retry.py.
\begin{table}[t]
\centering
\small
\setlength{\tabcolsep}{3.5pt}
\begin{tabular}{@{}lrrrr@{}}
\toprule
Schedule & \multicolumn{2}{c}{Delivered: immediate $\to$ age four} & \multicolumn{2}{c}{$\Delta$ bytes/delivery}\\
\cmidrule(lr){2-3}\cmidrule(lr){4-5}
 & $b=32$ & $b=64$ & $b=32$ & $b=64$\\
\midrule
Alternating & 688 $\to$ 971 & 718 $\to$ 999 & -11.41\% & -19.70\%\\
Lossy & 665 $\to$ 905 & 681 $\to$ 965 & +14.37\% & -3.99\%\\
Reordered & 671 $\to$ 784 & 700 $\to$ 837 & +3.94\% & -12.07\%\\
Duplicates & 688 $\to$ 971 & 718 $\to$ 999 & -11.41\% & -19.70\%\\
9:1 asymmetric & 517 $\to$ 648 & 720 $\to$ 857 & -32.12\% & -35.76\%\\
Offline then exchange & 613 $\to$ 879 & 653 $\to$ 935 & -22.57\% & -25.39\%\\
One-way & 32 $\to$ 32 & 64 $\to$ 64 & +0.09\% & +0.10\%\\
\bottomrule
\end{tabular}
\caption{Application-retry readiness at the same 2,048 directed opportunities. Each delivery column compares the freshly repeated immediate policy with minimum retry age four in the sender's own opportunities. Byte cost divides all emitted control, ACK, application and retry bytes by actual logical deliveries; positive changes indicate higher cost per delivery.}
\label{tab:retry-readiness}
\end{table}


\paragraph{Completed service and frame use.}
Age four increased logical deliveries by 16.84--43.39\% across the
twelve bidirectional or resumed budget--schedule pairs
(Table~\ref{tab:retry-readiness}). On the alternating schedule,
deliveries rose from 688 to 971 at $b=32$ (41.13\%) and from 718 to
999 at $b=64$ (39.14\%). Application retries fell from 604 to 16 and
from 579 to 8. Neither age-four run retried a message after its first
delivery. The immediate policy had done so 592 and 572 times.

The two age-four runs emitted 706,764 and 455,388 bytes, respectively.
Bytes per delivery fell by 11.41\% and 19.70\%. Each captured endpoint
still yielded 63 or 127 post-capture keys, out of 970 or 998 post-capture
first commitments. Physical \texttt{sync\_data} calls rose from 14,219
to 20,015 and from 14,579 to 20,255. Per delivery, those counts fell
from 20.67 to 20.61 and from 20.31 to 20.28.

\paragraph{Schedule-dependent accounting.}
On the lossy and reordered schedules at $b=32$, the endpoints delivered
36.09\% and 16.84\% more applications. They also spent 14.37\% and
3.94\% more bytes per delivery. A committed message can still be waiting
when a run ends. With age four, the two lossy runs committed/delivered
906/905 and 966/965 applications, and the $b=64$ asymmetric run reached
858/857. We count post-capture first commitments in the exposure
denominator. Under the offline prefix, the derived counts rose from
32 to 33 and from 64 to 65 as those denominators rose from 612 to 878
and from 652 to 934. The allocation histories changed within the same
$2b-1$ bound.

With traffic in one direction only, each policy admits and delivers
exactly 32 or 64 applications. The remaining requests stay queued and
both confirmed epochs stay at zero. Switching to age four adds one
4,746-byte control frame and saves one 141-byte retry, a net increase
of 4,605 bytes for the same service.

Appendix~\ref{app:retry-readiness} identifies the driver and the checks
that regenerate the table. The artifact records all 14 pairs, their
original counts and denominators, the identities of the 28 public traces
and an independent recount of retry ages and frame outcomes.
